% \iffalse meta-comment
%
% Copyright (C) 2017- by Yanshuo Chu <yanshuoc@gmail.com>
%
% This file may be distributed and/or modified under the
% conditions of the LaTeX Project Public License, either version 1.3a
% of this license or (at your option) any later version.
% The latest version of this license is in:
%
% http://www.latex-project.org/lppl.txt
%
% and version 1.3a or later is part of all distributions of LaTeX
% version 2004/10/01 or later.
%
% \fi
%
% ^^A ======================================================================
% ^^A hithesis 的 art / artplus 家族实现
% ^^A
% ^^A 开题、中期报告（hithesisart）与深圳博士中期（hithesisartplus）两个类的
% ^^A 全部代码，以及 art 类的配置文件 hithesisart.cfg。
% ^^A
% ^^A 三个类共用的文件头（\NeedsTeXFormat、\ProvidesClass 等）留在 hithesis.dtx，
% ^^A hithesis.ins 里每个 \file 先从那里取头，再从本文件取实现。
% ^^A ======================================================================
%
% ^^A ======================================================================
% ^^A Module artpluscls: artplus 类主体（深圳博士中期报告专用）
% ^^A ======================================================================
% ^^A ======================================================================
% ^^A Module artplus-options: 文档类选项的声明与后处理（type/campus/fontset 等）（artplus）
% ^^A ======================================================================
%    \begin{macrocode}
%<*artplus-options>
%    \end{macrocode}
%
% 处理文档类选项
%    \begin{macrocode}

%% 处理文档类选项
\RequirePackage{kvoptions}
\SetupKeyvalOptions{
  family=hit,
  prefix=hit@,
  setkeys=\kvsetkeys
}

\newif\ifhit@bachelor
\newif\ifhit@master
\newif\ifhit@doctor
\define@key{hit}{type}{%
  \hit@bachelorfalse
  \hit@masterfalse
  \hit@doctorfalse
  \expandafter\csname hit@#1true\endcsname
}

\newif\ifhit@shenzhen
\newif\ifhit@weihai
\newif\ifhit@harbin
\define@key{hit}{campus}{%
  \hit@shenzhenfalse
  \hit@weihaifalse
  \hit@harbinfalse
  \expandafter\csname hit@#1true\endcsname
}
\ifhit@weihai\relax\else
  \ifhit@shenzhen\relax\else
    \hit@harbintrue
  \fi
\fi

\newif\ifhit@opening
\newif\ifhit@midterm
\define@key{hit}{stage}{%
  \hit@openingfalse
  \hit@midtermfalse
  \expandafter\csname hit@#1true\endcsname
}
%    \end{macrocode}
% \changes{v3.2a}{0000/00/00}{设置参考文献列表中的作者是否全大写和最大显示数量}
% \changes{v3.2a}{0000/00/00}{设置参考文献列表标号的对齐方式}
%    \begin{macrocode}

\ExplSyntaxOn
\bool_new:N \g_bib_authorupper_bool
\bool_set_false:N \g_bib_authorupper_bool
\define@key{hit}{bibauthorupper}{%
  \str_if_eq:nnT {#1} {full} {
    \bool_set_true:N \g_bib_authorupper_bool
  }
}

\int_new:N \g_bib_maxauthor_int
\int_set:Nn \g_bib_maxauthor_int {3}
\define@key{hit}{bibmaxauthor}{%
  \int_set:Nn \g_bib_maxauthor_int {#1}
  \int_if_zero:nT {\g_bib_maxauthor_int} {
    \int_gdecr:N \g_bib_maxauthor_int
  }
}

\str_new:N \g_bib_labelalign_str
\define@key{hit}{biblabelalign}{%
  \str_set:Nn \g_bib_labelalign_str {#1}
}
\ExplSyntaxOff

\DeclareStringOption{fontset}
\DeclareDefaultOption{\PassOptionsToClass{\CurrentOption}{ctexart}}
\PassOptionsToPackage{no-math}{fontspec}

\ProcessKeyvalOptions*

\ifhit@bachelor\relax\else
  \ifhit@master\relax\else
    \ifhit@doctor\relax\else
      \ClassError{hithesisartplus}{%
        \MessageBreak Please specify thesis type in option:
        \MessageBreak type=[bachelor | master | doctor]
      }
    \fi
  \fi
\fi

\ifhit@opening\relax\else
  \ifhit@midterm\relax\else
    \ClassError{hithesisartplus}{%
      \MessageBreak Please specify stage in option:
      \MessageBreak stage=<opening|midterm>
    }
    \fi
\fi

\ifhit@doctor
  \ifhit@midterm
    \ifhit@shenzhen
      \relax
    \else
      \ClassError{hithesisartplus}{%
        \MessageBreak This document class only support midterm report for doctor
        in shenzhen campus
      }
    \fi
  \else
    \ClassError{hithesisartplus}{%
      \MessageBreak This document class only support midterm report for doctor
      in shenzhen campus
    }
  \fi
\else
  \ClassError{hithesisartplus}{%
    \MessageBreak This document class only support midterm report for doctor
    in shenzhen campus
  }
\fi
%    \end{macrocode}
%
% \changes{v3.2a}{0000/00/00}{选项后处理（AutoFakeBold 与 fontset 透传）从 artpluscls 移入 artplus-options 模块}
%    \begin{macrocode}
\RequirePackage{ifthen}
\ifthenelse
{\equal{\hit@fontset}{}}%
{%
  \PassOptionsToPackage{AutoFakeBold=2}{xeCJK}
}%
{%
  \ifthenelse
  {\equal{\hit@fontset}{siyuan}}%
  {\relax}%
  {%
    \PassOptionsToPackage{AutoFakeBold=2}{xeCJK}
  }%
  \PassOptionsToClass{fontset=\hit@fontset}{ctexart}
}%
%</artplus-options>
%<*artpluscls-load>
\LoadClass[a4paper,UTF8,zihao=-4,scheme=plain]{ctexart}
%    \end{macrocode}
% \changes{v3.2a}{0000/00/00}{模块内容改由 \file{hithesis.ins} 按运行期顺序直接拼入 cls，不再生成 \file{modules/*.sty}}
%    \begin{macrocode}
%    \end{macrocode}
%
% \changes{v3.2a}{0000/00/00}{LoadClass 之后的所有标准包加载、配置整体打包至 artplus-deps 模块；artpluscls 仅保留 LoadClass}
%    \begin{macrocode}
%</artpluscls-load>
%    \end{macrocode}
% ^^A ======================================================================
% ^^A Module artplus-deps: 第三方宏包集中加载与基础配置（含 hyperlink/geometry 内联）（artplus）
% ^^A ======================================================================
%    \begin{macrocode}
%<*artplus-deps-a>
%    \end{macrocode}
%
% 各种宏包
%    \begin{macrocode}

%% 各种宏包
\RequirePackage{etoolbox}
\RequirePackage{ifxetex}
\ifxetex
\else
  \ClassError{hithesis}{Please use: \MessageBreak xelatex}
\fi
\RequirePackage{xparse}

\RequirePackage{amsmath}
\RequirePackage[amsmath,thmmarks,hyperref]{ntheorem}
\RequirePackage{amssymb}
%    \end{macrocode}
%
% 删除 \pkg{newtxtext}，由于最新一次更新 (\texttt{revision 62369}) 中与 \pkg{ctex} 宏包冲突，且无法解决，将正文字体替换为 TG TermesX 与 TG Heros.
% \changes{v3.0.18}{2022/05/02}{删除 \pkg{newtxtext}，由于最新一次更新 (\texttt{revision 62369}) 中与 \pkg{ctex} 宏包冲突，且无法解决，将正文字体替换为 TG Termes 与 TG Heros.}
% \changes{v3.0.19}{2022/05/05}{将 TG Termes 字体替换为 TG TermesX 字体，与 \pkg{newtxmath} 宏包更适配，删除 \option{Scale} 选项。}
% \changes{v3.2a}{0000/00/00}{英文字体三件套设置从 artpluscls 拆出至 artplus-fonts 模块}
%    \begin{macrocode}
%</artplus-deps-a>
%    \end{macrocode}
% ^^A ======================================================================
% ^^A Module artplus-fonts: 正文字号、中英文主字体与字号宏命令（\wuhao 等）（artplus）
% ^^A ======================================================================
%    \begin{macrocode}
%<*artplus-fonts>
\setmainfont{TeXGyreTermesX}[
  UprightFont = *-Regular ,
  BoldFont = *-Bold ,
  ItalicFont = *-Italic ,
  BoldItalicFont = *-BoldItalic ,
  Extension = .otf ,
  ]
\setsansfont{texgyreheros}[
  UprightFont = *-regular ,
  BoldFont = *-bold ,
  ItalicFont = *-italic ,
  BoldItalicFont = *-bolditalic ,
  Extension = .otf ,
  ]
%    \end{macrocode}
%
% 删除 \pkg{courier}，将 mono 字体替换为 TG Cursor，至此将正文字体全面替换为 OpenType.
% \changes{v3.0.19}{2022/05/05}{删除 \pkg{courier}，将 mono 字体替换为 TG Cursor，至此将正文字体全面替换为 OpenType.}
%    \begin{macrocode}
\setmonofont{texgyrecursor}[
  UprightFont = *-regular ,
  BoldFont = *-bold ,
  ItalicFont = *-italic ,
  BoldItalicFont = *-bolditalic ,
  Extension = .otf ,
]

%% 添加字号、行距等设置
\renewcommand\normalsize{%
  \@setfontsize\normalsize{12bp}{15bp}%
  \abovedisplayskip=6pt
  \abovedisplayshortskip=6pt
  \belowdisplayskip=\abovedisplayskip
  \belowdisplayshortskip=\abovedisplayshortskip
}
\def\hit@def@fontsize#1#2{%
  \expandafter\newcommand\csname #1\endcsname[1][1.3]{%
    \fontsize{#2}{##1\dimexpr #2}\selectfont
  }
}
\hit@def@fontsize{dachu}{58bp}
\hit@def@fontsize{chuhao}{42bp}
\hit@def@fontsize{xiaochu}{36bp}
\hit@def@fontsize{yihao}{26bp}
\hit@def@fontsize{xiaoyi}{24bp}
\hit@def@fontsize{erhao}{22bp}
\hit@def@fontsize{xiaoer}{18bp}
\hit@def@fontsize{sanhao}{16bp}
\hit@def@fontsize{xiaosan}{15bp}
\hit@def@fontsize{sihao}{14bp}
\hit@def@fontsize{banxiaosi}{13bp}
\hit@def@fontsize{xiaosi}{12bp}
\hit@def@fontsize{dawu}{11bp}
\hit@def@fontsize{wuhao}{10.5bp}
\hit@def@fontsize{xiaowu}{9bp}
\hit@def@fontsize{liuhao}{7.5bp}
\hit@def@fontsize{xiaoliu}{6.5bp}
\hit@def@fontsize{qihao}{5.5bp}
\hit@def@fontsize{bahao}{5bp}
%</artplus-fonts>
%<*artplus-deps-a>

\RequirePackage{graphicx}
\RequirePackage{pdfpages}
\includepdfset{fitpaper=true}
\RequirePackage{enumitem}       %使用enumitem宏包,改变列表项的格式
\RequirePackage{environ}
\RequirePackage[perpage,hang]{footmisc}
%    \end{macrocode}
%
% \pkg{xeCJK} 包不再自动加载 \pkg{xeCJKfntef}
% \changes{v3.0.17}{2022/02/23}{\pkg{xeCJK} 包不再自动加载 \pkg{xeCJKfntef}}
%    \begin{macrocode}
\RequirePackage{xeCJKfntef}
%    \end{macrocode}
%
% \pkg{newtx} 更新至 v1.7 版本后与 \pkg{ctex} 宏包冲突, 需要手动引入 \option{nofontspec} 选项. 判断方法为是否存在 v1.7 更新后的新宏包 \pkg{newtx.sty}
% \changes{v3.0.17}{2022/02/23}{\pkg{newtx} 更新至 v1.7 版本后与 \pkg{ctex} 宏包冲突, 需要手动引入 \option{nofontspec} 选项. 判断方法为是否存在 v1.7 更新后的新宏包 \pkg{newtx.sty}}
%    \begin{macrocode}

\RequirePackage{longtable}
\RequirePackage{booktabs}
%    \end{macrocode}
%
% 重定义 \pkg{booktabs} 宏包的默认线宽，绘制三线表时不再需要显式设置，以实现更好的内容与格式分离
% \changes{v3.0.20}{2022/05/06}{重定义 \pkg{booktabs} 宏包的默认线宽，绘制三线表时不再需要显式设置，以实现更好的内容与格式分离}
%    \begin{macrocode}
\heavyrulewidth=1.5pt
\lightrulewidth=1pt
\cmidrulewidth=0.5pt

%    \end{macrocode}
%
% \changes{v3.2a}{0000/00/00}{natbib/hyperref/url 配置从 artplus-deps 拆出至 artplus-hyperlink 模块（对齐 book-hyperlink/-art）}
%    \begin{macrocode}
%</artplus-deps-a>
%    \end{macrocode}
% ^^A ======================================================================
% ^^A Module artplus-hyperlink: hyperref/url 配置（PDF 书签、链接颜色）（artplus）
% ^^A ======================================================================
%    \begin{macrocode}
%<*artplus-hyperlink>
\RequirePackage{hyperref}
\hypersetup{%
  CJKbookmarks=true,
  linktoc=all,
  bookmarksnumbered=true,
  bookmarksopen=true,
  bookmarksopenlevel=1,
  breaklinks=true,
  colorlinks=false,
  plainpages=false,
  pdfborder=0 0 0
}
\urlstyle{same}
%</artplus-hyperlink>
%<*artplus-deps-b>

%    \end{macrocode}
%
% \changes{v3.2a}{0000/00/00}{geometry 配置从 artplus-deps 拆出至 artplus-geometry 模块（对齐 book-geometry/-art）}
%    \begin{macrocode}
%</artplus-deps-b>
%    \end{macrocode}
% ^^A ======================================================================
% ^^A Module artplus-geometry: 版芯尺寸与页边距（artplus）
% ^^A ======================================================================
%    \begin{macrocode}
%<*artplus-geometry>
\RequirePackage{geometry}
\geometry{%根据PlutoThesis 原版定义而来
  a4paper, % 210 * 297mm
  hcentering,
  ignoreall,
  nomarginpar,
  centering,
  text={150true mm,236true mm},
  left=30true mm,
  head=5true mm,
  headsep=2true mm,
  footskip=0true mm,
  foot=5.2true mm
}
%</artplus-geometry>
%<*artplus-deps-c>

\RequirePackage{tikzpagenodes}
\RequirePackage{eso-pic}
\RequirePackage{tikz}
\usetikzlibrary{calc}

\RequirePackage{fancyhdr}
\RequirePackage{tabularx}
\RequirePackage{varwidth}
\RequirePackage{changepage}
\RequirePackage{multicol}
\RequirePackage[below]{placeins}%允许上一个section的浮动图形出现在下一个section的开始部分,还提供\FloatBarrier命令,使所有未处理的浮动图形立即被处理
\RequirePackage{flafter}       % 使得所有浮动体不能被放置在其浮动环境之前，以免浮动体在引述它的文本之前出现.
\RequirePackage{multirow}       %使用Multirow宏包，使得表格可以合并多个row格
\RequirePackage{subfigure}%支持子图 %centerlast 设置最后一行是否居中
\RequirePackage[subfigure]{ccaption} %支持双语标题
%    \end{macrocode}
% % \changes{v3.0.21}{2022/05/11}{删除 \pkg{xltxtra} 宏包, 来实现规范样式的上标。 \issue[125]}
%    \begin{macrocode}

\RequirePackage{lipsum}%random text
%    \end{macrocode}
%
% 添加字号、行距等设置
%    \begin{macrocode}
%</artplus-deps-c>
%<*artplus-deps-c>
%    \end{macrocode}
%
% 进行字号、行距等设置
% \changes{v3.2a}{0000/00/00}{章节标题格式从 artpluscls 拆出至 artplus-chapter 模块}
%    \begin{macrocode}
%</artplus-deps-c>
%    \end{macrocode}
% ^^A ======================================================================
% ^^A Module artplus-chapter: 章节标题格式（字号、缩进、间距、编号样式）（artplus）
% ^^A ======================================================================
%    \begin{macrocode}
%<*artplus-chapter>
\RequirePackage{xparse}
\RequirePackage{tikz}

%% 进行字号、行距等设置
\ctexset{%
  section={
    afterindent=true,
    beforeskip={7.5bp},%上下空0.5行
    afterskip={7.5bp},
    format={\normalsize},
    aftername=\enspace,
    fixskip=true,
    break={},
  },
  subsection={
    afterindent=true,
    beforeskip={7bp},
    afterskip={7bp},
    format={\normalsize},
    aftername=\enspace,
    fixskip=true,
    break={},
  },
  subsubsection={
    afterindent=true,
    beforeskip={3bp},
    afterskip={3bp},
    format={\normalsize},
    aftername=\enspace,
    fixskip=true,
    break={},
  },
  paragraph/afterindent=true,
  subparagraph/afterindent=true
}

\let\hit@section\section
\RenewDocumentCommand\section{s o m s}{%
  \IfBooleanF{#4}{\noindent\tikz[overlay] \draw (-5mm,0mm) -- (155mm,0mm);}%
  \IfBooleanTF{#1}%
  {% if \section*
    \IfNoValueTF{#2}%
    {\hit@section*{#3}}%
    {\hit@section*[#2]{#3}}%
  }%
  {% if \section
    \IfNoValueTF{#2}%
    {\hit@section{#3}}%
    {\hit@section[#2]{#3}}%
  }
}
%</artplus-chapter>
%<*artplus-deps-c>
%    \end{macrocode}
%
% 中文封面
%    \begin{macrocode}
%</artplus-deps-c>
%    \end{macrocode}
% ^^A ======================================================================
% ^^A Module artplus-meta: 元信息字段（\hitsetup 接受的 key）（artplus）
% ^^A ======================================================================
%    \begin{macrocode}
%<*artplus-meta>

%% 中文封面
\def\hit@def@term#1{%
  \define@key{hit}{#1}{\csname #1\endcsname{##1}}
  \expandafter\gdef\csname #1\endcsname##1{%
    \expandafter\gdef\csname hit@#1\endcsname{##1}
  }
  \csname #1\endcsname{}
}
%    \end{macrocode}
% \changes{v3.1d}{2025/03/03}{取消了深圳博士中期——即reportplus——中未使用到的参数ctitlecover}
%    \begin{macrocode}
\hit@def@term{ctitleone}
\hit@def@term{ctitletwo}
\hit@def@term{csubject}
\hit@def@term{cauthor}
\hit@def@term{cstudentid}
\hit@def@term{cclassid}
\hit@def@term{caffil}
\hit@def@term{csupervisor}
\hit@def@term{cdate}
\hit@def@term{cenrolldate}

\def\hit@parse@keywords#1{
  \define@key{hit}{#1}{\csname #1\endcsname{##1}}
  \expandafter\gdef\csname hit@#1\endcsname{}
  \expandafter\gdef\csname #1\endcsname##1{
    \@for\reserved@a:=##1\do{
      \expandafter\ifx\csname hit@#1\endcsname\@empty\else
      \expandafter\g@addto@macro\csname hit@#1\endcsname{%
        \ignorespaces\csname hit@#1@separator\endcsname}
      \fi
      \expandafter\expandafter\expandafter\g@addto@macro%
      \expandafter\csname hit@#1\expandafter\endcsname\expandafter{\reserved@a}
    }
  }
}

\def\hitsetup{\kvsetkeys{hit}}
%</artplus-meta>
%    \end{macrocode}
% ^^A ======================================================================
% ^^A Module artplus-pagestyle: 页眉页脚样式（artplus）
% ^^A ======================================================================
% \changes{v3.2a}{0000/00/00}{页眉页脚样式从 artplus-frontmatter 的 \cs{makecover} 内联提取到独立的 artplus-pagestyle 模块（对齐 book-pagestyle/-art）}
%    \begin{macrocode}
%<*artplus-pagestyle>
\RequirePackage{fancyhdr}
\fancypagestyle{hit@artplus@main}{%
  \fancyhf{}%
  \renewcommand{\headrulewidth}{0pt}%
  \renewcommand{\footrulewidth}{0pt}%
  \fancyfoot[C]{\xiaowu-~\thepage~-}%
}
%</artplus-pagestyle>
%    \end{macrocode}
% ^^A ======================================================================
% ^^A Module artplus-frontmatter: 封面与扉页（中英、本科、博士后等多版本）（artplus）
% ^^A ======================================================================
%    \begin{macrocode}
%<*artplus-frontmatter>

\newcommand{\hit@report@titlepage@graduate}{
  \ifthenelse
  {\equal{\hit@fontset}{siyuan}}%
  {\xiaosi[1]\vspace*{0.65em}}%
  {\xiaosi[1]\textcolor[rgb]{1,1,1}{\songti{\hit@hi}}}%
  \vspace*{10mm}
  \begin{center}
    \songti\erhao\textbf{\hit@cschoolname\ifhit@shenzhen\hit@shenzhencampus\fi}
  \end{center}
  \begin{center}
    \songti\erhao\textbf{\hit@cthesismidtermcheck}
  \end{center}
  \vspace{66bp}
  \parbox[b][3cm][t]{\textwidth}{
    \begin{center}\songti\sihao
      \renewcommand{\arraystretch}{1.67}
      \begin{tabular}{l@{\ \ }c}
        \textbf{\hit@graduate@studenttitle} & \textbf{\hit@cauthor}\\
        \hline
        \textbf{\hit@graduate@enrolldate} & \textbf{\hit@cenrolldate}\\
        \hline
        \textbf{\hit@graduate@thesistitle} & \textbf{\hit@ctitleone}\\
        \hline
        &\textbf{\hit@ctitletwo}\\
        \hline
        \textbf{\hit@graduate@cafflimajor} & \textbf{\hit@csubject}\\
        \hline
        \textbf{\hit@graduate@datetitle} & \textbf{\hit@cdate}\\
        \hline
        \textbf{\hit@graduate@supervisor} & \textbf{\hit@csupervisor}\\
        \hline
      \end{tabular}\renewcommand{\arraystretch}{1}
    \end{center}
  }
}

\def\makecover{
  \begin{titlepage}
    \hit@report@titlepage@graduate
    \clearpage
    \hit@shenzhen@doctor@midterm@note
  \end{titlepage}
  \AddToShipoutPicture{%
    \begin{tikzpicture}[remember picture,overlay]
      \draw ($(current page text area.south west)+(-5mm,0mm)$) rectangle ($(current page text area.north east)+(5mm,2mm)$);
    \end{tikzpicture}
  }
  \pagestyle{hit@artplus@main}
}
%</artplus-frontmatter>
%<*artplus-deps-c>
%    \end{macrocode}
%
% 引用和参考文献格式
% \changes{v3.2a}{0000/00/00}{设置参考文献列表标号的对齐方式}
% \changes{v3.2a}{0000/00/00}{参考文献格式从 artpluscls 拆出至 artplus-bib 模块}
%    \begin{macrocode}
%</artplus-deps-c>
%    \end{macrocode}
% ^^A ======================================================================
% ^^A Module artplus-bib: 参考文献加载与 natbib/gbt7714 设置（artplus）
% ^^A ======================================================================
%    \begin{macrocode}
%<*artplus-bib>
\RequirePackage[sort&compress,numbers]{natbib}
\RequirePackage{xparse}

%% 引用和参考文献格式
\newcommand\bibstyle@numerical{\bibpunct{[}{]}{,}{s}{,}{\textsuperscript{,}}}
\newcommand\bibstyle@inline{\bibpunct{[}{]}{,}{n}{,}{\hit@inline@sep}}

\citestyle{numerical}
\DeclareRobustCommand\inlinecite{\@inlinecite}
\def\@inlinecite#1{\begingroup\citestyle{inline}\let\@cite\NAT@citenum\citep{#1}\endgroup}
\let\onlinecite\inlinecite

\ExplSyntaxOn
\renewenvironment{thebibliography}[1]{%
  \list{\@biblabel{\@arabic\c@enumiv}}{%
    \str_if_eq:VnT \g_bib_labelalign_str {left} {
      \RenewDocumentCommand\makelabel{m}{##1\hfill}
    }
    \settowidth{\labelwidth}{\@biblabel{#1}}
    \setlength{\labelsep}{0.5em}
    \setlength{\itemindent}{0pt}
    \setlength{\leftmargin}{\labelsep+\labelwidth}
    \addtolength{\itemsep}{-0.8em}
    \usecounter{enumiv}%
    \let\p@enumiv\@empty
  \renewcommand\theenumiv{\@arabic\c@enumiv}}%
  \sloppy\frenchspacing
  \flushbottom
  \clubpenalty0
  \@clubpenalty \clubpenalty
  \widowpenalty0%
  \interlinepenalty-50%
  \sfcode`\.\@m
}{}
%    \end{macrocode}
% \changes{v3.2a}{0000/00/00}{为上面的环境补充了一对空大括号}
% \changes{v3.2a}{0000/00/00}{控制参考文献列表中作者是否全大写和最大显示数量}
%    \begin{macrocode}
%% 用以控制参考文献列表中单个作者的格式
\exp_args:NNe \cs_new:Nn \g_bib_author_cs:n {
    \bool_if:NTF \g_bib_authorupper_bool {
        \exp_not:N \text_uppercase:n {#1}
    } {
        {#1}
    }
}
%% 用以输出参考文献列表中的作者
\NewDocumentCommand{\bibauthor}{m}{
  \int_zero:N \l_tmpa_int
  \exp_args:Ne \tl_map_inline:nn {
    \exp_args:Ne \tl_tail:n {#1}
  } {
    \int_compare:nNnT
      \l_tmpa_int = \g_bib_maxauthor_int {
      \prg_break:n {\tl_head:N {#1}}
    }
    \int_incr:N \l_tmpa_int
    \g_bib_author_cs:n {##1}
  }
  \prg_break_point:
}
\ExplSyntaxOff
%</artplus-bib>
%<*artpluscls-tail>
%    \end{macrocode}
% 杂项
% \changes{v3.2a}{0000/00/00}{misc-artplus 模块合并为 artpluscls 内联钩子}
%    \begin{macrocode}
\AtEndOfClass{\input{hithesisart.cfg}}
\AtEndOfClass{\sloppy}
%</artpluscls-tail>
%    \end{macrocode}
%
% 添加artcls
% \changes{v3.0.0}{2020/05/21}{添加artcls}
%    \begin{macrocode}
% %%%%%%%%%%%%%%%%%%%%%%%%%%%%%%%%%%%%%%%%%%%%%%%%%%%%%%%%%%%%%%%%%%%%%%%%%%%%%%
% artcls
% %%%%%%%%%%%%%%%%%%%%%%%%%%%%%%%%%%%%%%%%%%%%%%%%%%%%%%%%%%%%%%%%%%%%%%%%%%%%%%
%    \end{macrocode}
% ^^A ======================================================================
% ^^A Module artcls: art 类主体（开题/中期报告）
% ^^A ======================================================================
% ^^A ======================================================================
% ^^A Module art-options: 文档类选项的声明与后处理（type/campus/fontset 等）（art）
% ^^A ======================================================================
%    \begin{macrocode}
%<*art-options>
%    \end{macrocode}
%
% 处理文档类选项
%    \begin{macrocode}

%% 处理文档类选项
\RequirePackage{kvoptions}
\SetupKeyvalOptions{
  family=hit,
  prefix=hit@,
  setkeys=\kvsetkeys
}

\newif\ifhit@bachelor
\newif\ifhit@master
\newif\ifhit@doctor
\define@key{hit}{type}{%
  \hit@bachelorfalse
  \hit@masterfalse
  \hit@doctorfalse
  \expandafter\csname hit@#1true\endcsname
}

\newif\ifhit@shenzhen
\newif\ifhit@weihai
\newif\ifhit@harbin
\define@key{hit}{campus}{%
  \hit@shenzhenfalse
  \hit@weihaifalse
  \hit@harbinfalse
  \expandafter\csname hit@#1true\endcsname
}
\ifhit@weihai\relax\else
  \ifhit@shenzhen\relax\else
    \hit@harbintrue
  \fi
\fi

\newif\ifhit@opening
\newif\ifhit@midterm
\define@key{hit}{stage}{%
  \hit@openingfalse
  \hit@midtermfalse
  \expandafter\csname hit@#1true\endcsname
}
%    \end{macrocode}
% \changes{v3.2a}{0000/00/00}{设置参考文献列表中的作者是否全大写和最大显示数量}
% \changes{v3.2a}{0000/00/00}{设置参考文献列表标号的对齐方式}
%    \begin{macrocode}

\ExplSyntaxOn
\bool_new:N \g_bib_authorupper_bool
\bool_set_false:N \g_bib_authorupper_bool
\define@key{hit}{bibauthorupper}{%
  \str_if_eq:nnT {#1} {full} {
    \bool_set_true:N \g_bib_authorupper_bool
  }
}

\int_new:N \g_bib_maxauthor_int
\int_set:Nn \g_bib_maxauthor_int {3}
\define@key{hit}{bibmaxauthor}{%
  \int_set:Nn \g_bib_maxauthor_int {#1}
  \int_if_zero:nT {\g_bib_maxauthor_int} {
    \int_gdecr:N \g_bib_maxauthor_int
  }
}

\str_new:N \g_bib_labelalign_str
\define@key{hit}{biblabelalign}{%
  \str_set:Nn \g_bib_labelalign_str {#1}
}
\ExplSyntaxOff

\DeclareBoolOption[true]{raggedbottom}
\DeclareBoolOption[false]{pifootnote}
\DeclareBoolOption[false]{debug}
\DeclareBoolOption[true]{toc}
\DeclareBoolOption[true]{newtxmath}
%    \end{macrocode}

% \changes{v3.1c}{2023/04/16}{增加用于打印的选项 \meta{print}.\issue[138]}
%    \begin{macrocode}
\DeclareBoolOption[false]{print}


\DeclareStringOption{fontset}
\DeclareDefaultOption{\PassOptionsToClass{\CurrentOption}{ctexart}}
\PassOptionsToPackage{no-math}{fontspec}

\ProcessKeyvalOptions*

\ifhit@bachelor\relax\else
  \ifhit@master\relax\else
    \ifhit@doctor\relax\else
    \ClassError{hithesisart}{%
      \MessageBreak Please specify thesis type in option:
      \MessageBreak type=[bachelor | master | doctor]
    }
    \fi
  \fi
\fi

\ifhit@opening\relax\else
  \ifhit@midterm\relax\else
    \ClassError{hithesisart}{%
      \MessageBreak Please specify stage in option:
      \MessageBreak stage=<opening|midterm>
    }
  \fi
\fi

\ifhit@doctor
  \ifhit@midterm
    \ifhit@shenzhen
      \ClassError{hithesisart}{%
        \MessageBreak This document class does not support midterm report for doctor
        in shenzhen campus.
        \MessageBreak please use \string\documentclass{hithesisartplus}
      }
    \fi
  \fi
\fi
%    \end{macrocode}
%
% \changes{v3.2a}{0000/00/00}{选项后处理（AutoFakeBold 与 fontset 透传）从 artcls 移入 art-options 模块}
%    \begin{macrocode}
\RequirePackage{ifthen}
\ifthenelse
{\equal{\hit@fontset}{}}%
{%
  \PassOptionsToPackage{AutoFakeBold=2}{xeCJK}
}%
{%
  \ifthenelse
  {\equal{\hit@fontset}{siyuan}}%
  {\relax}%
  {%
    \PassOptionsToPackage{AutoFakeBold=2}{xeCJK}
  }%
  \PassOptionsToClass{fontset=\hit@fontset}{ctexart}
}%
%</art-options>
%<*artcls-load>
\LoadClass[a4paper,UTF8,zihao=-4,scheme=plain]{ctexart}
%    \end{macrocode}
% \changes{v3.2a}{0000/00/00}{模块内容改由 \file{hithesis.ins} 按运行期顺序直接拼入 cls，不再生成 \file{modules/*.sty}}
%    \begin{macrocode}
%    \end{macrocode}
%
% \changes{v3.2a}{0000/00/00}{LoadClass 之后的所有标准包加载、配置整体打包至 art-deps 模块；artcls 仅保留 LoadClass}
% \changes{v3.2a}{0000/00/00}{目录排版从 art-frontmatter 拆出至 art-toc 模块（对齐 book 家族的 book-toc）；排在 art-deps 之后以避免 enumitem 覆盖 cft* 类长度}
%    \begin{macrocode}
%</artcls-load>
%    \end{macrocode}
% ^^A ======================================================================
% ^^A Module art-deps: 第三方宏包集中加载与基础配置（含 hyperlink/geometry 内联）（art）
% ^^A ======================================================================
%    \begin{macrocode}
%<*art-deps-a>
%    \end{macrocode}
%
% 各种宏包
%    \begin{macrocode}

%% 各种宏包
\RequirePackage{etoolbox}
\RequirePackage{ifxetex}
\ExplSyntaxOn
%    \end{macrocode}
%
% \changes{v3.2a}{0000/00/00}{art-deps-a 的引擎检查改用 \cs{legacy_if:nF}。报错
%   文本里的空格在 \cs{ExplSyntaxOn} 下会被吞，所以写成 \texttt{\~{}}。}
%    \begin{macrocode}
\legacy_if:nF { xetex }
  { \ClassError { hithesis } { Please~use:~\MessageBreak~xelatex } { } }
\ExplSyntaxOff
\RequirePackage{xparse}

\RequirePackage{amsmath}
\RequirePackage[amsmath,thmmarks,hyperref]{ntheorem}
\RequirePackage{amssymb}
%    \end{macrocode}
%
% 删除 \pkg{newtxtext}，由于最新一次更新 (\texttt{revision 62369}) 中与 \pkg{ctex} 宏包冲突，且无法解决，将正文字体替换为 TG TermesX 与 TG Heros.
% \changes{v3.0.18}{2022/05/02}{删除 \pkg{newtxtext}，由于最新一次更新 (\texttt{revision 62369}) 中与 \pkg{ctex} 宏包冲突，且无法解决，将正文字体替换为 TG Termes 与 TG Heros.}
% \changes{v3.0.19}{2022/05/05}{将 TG Termes 字体替换为 TG TermesX 字体，与 \pkg{newtxmath} 宏包更适配，删除 \option{Scale} 选项。}
% \changes{v3.2a}{0000/00/00}{英文字体三件套设置从 artcls 拆出至 art-fonts 模块}
%    \begin{macrocode}
%</art-deps-a>
%    \end{macrocode}
% ^^A ======================================================================
% ^^A Module art-fonts: 正文字号、中英文主字体与字号宏命令（\wuhao 等）（art）
% ^^A ======================================================================
%    \begin{macrocode}
%<*art-fonts>
\setmainfont{TeXGyreTermesX}[
  UprightFont = *-Regular ,
  BoldFont = *-Bold ,
  ItalicFont = *-Italic ,
  BoldItalicFont = *-BoldItalic ,
  Extension = .otf ,
  ]
\setsansfont{texgyreheros}[
  UprightFont = *-regular ,
  BoldFont = *-bold ,
  ItalicFont = *-italic ,
  BoldItalicFont = *-bolditalic ,
  Extension = .otf ,
  ]
%    \end{macrocode}
%
% 删除 \pkg{courier}，将 mono 字体替换为 TG Cursor，至此将正文字体全面替换为 OpenType.
% \changes{v3.0.19}{2022/05/05}{删除 \pkg{courier}，将 mono 字体替换为 TG Cursor，至此将正文字体全面替换为 OpenType.}
%    \begin{macrocode}
\setmonofont{texgyrecursor}[
  UprightFont = *-regular ,
  BoldFont = *-bold ,
  ItalicFont = *-italic ,
  BoldItalicFont = *-bolditalic ,
  Extension = .otf ,
]

%% 添加字号、行距等设置
\renewcommand\normalsize{%
  \@setfontsize\normalsize{12bp}{20.50398bp}%
  \abovedisplayskip=8pt
  \abovedisplayshortskip=8pt
  \belowdisplayskip=\abovedisplayskip
  \belowdisplayshortskip=\abovedisplayshortskip
}
\def\hit@def@fontsize#1#2{%
  \expandafter\newcommand\csname #1\endcsname[1][1.3]{%
    \fontsize{#2}{##1\dimexpr #2}\selectfont
  }
}
\hit@def@fontsize{dachu}{58bp}
\hit@def@fontsize{chuhao}{42bp}
\hit@def@fontsize{xiaochu}{36bp}
\hit@def@fontsize{yihao}{26bp}
\hit@def@fontsize{xiaoyi}{24bp}
\hit@def@fontsize{erhao}{22bp}
\hit@def@fontsize{xiaoer}{18bp}
\hit@def@fontsize{sanhao}{16bp}
\hit@def@fontsize{xiaosan}{15bp}
\hit@def@fontsize{sihao}{14bp}
\hit@def@fontsize{banxiaosi}{13bp}
\hit@def@fontsize{xiaosi}{12bp}
\hit@def@fontsize{dawu}{11bp}
\hit@def@fontsize{wuhao}{10.5bp}
\hit@def@fontsize{xiaowu}{9bp}
\hit@def@fontsize{liuhao}{7.5bp}
\hit@def@fontsize{xiaoliu}{6.5bp}
\hit@def@fontsize{qihao}{5.5bp}
\hit@def@fontsize{bahao}{5bp}
%</art-fonts>
%<*art-deps-a>
%    \end{macrocode}
%
% 单独使用 \pkg{newtxmath} 宏包改变数学字体时，会使 \texttt{operators}，\cs{mathrm}，\cs{mathit}，\cs{mathbf} 命令使用 CM 字体，根据 \href{https://github.com/sjtug/SJTUThesis/blob/b164d0f5b3087ed86c1038b5a85014b47a92c6f2/texmf/tex/latex/sjtuthesis/fd/sjtu-math-font-termes.def#L52-L65}{SJTUThesis 的字体设置} 进行调整
% \changes{v3.0.19}{2022/05/05}{对 \pkg{newtxmath} 宏包打补丁，修复 \texttt{operators}，\cs{mathrm}，\cs{mathit}，\cs{mathbf} 命令的字体问题。}
%    \begin{macrocode}
\ifhit@newtxmath
  \let\oldencodingdefault\encodingdefault
  \let\oldrmdefault\rmdefault
  \let\oldsfdefault\sfdefault
  \let\oldttdefault\ttdefault
%    \end{macrocode}
%
% \changes{v3.1b}{2022/06/06}{直接修改编码，而不引入新的宏包。}
%    \begin{macrocode}
  \def\encodingdefault{T1}
  \renewcommand{\rmdefault}{ntxtlf}
  \renewcommand{\sfdefault}{qhv}
  \renewcommand{\ttdefault}{ntxtt}
  \RequirePackage{newtxmath}
  \let\encodingdefault\oldencodingdefault
  \let\rmdefault\oldrmdefault
  \let\sfdefault\oldsfdefault
  \let\ttdefault\oldttdefault
\fi

\RequirePackage{graphicx}
\RequirePackage{pdfpages}
\includepdfset{fitpaper=true}
\RequirePackage{enumitem}       %使用enumitem宏包,改变列表项的格式
\RequirePackage{environ}
\ifhit@raggedbottom
  \RequirePackage[bottom,perpage,hang]{footmisc}
  \raggedbottom
\else
  \RequirePackage[perpage,hang]{footmisc}
\fi
\ifhit@pifootnote
  \RequirePackage{pifont}
\fi
%    \end{macrocode}
%
% \pkg{xeCJK} 包不再自动加载 \pkg{xeCJKfntef}
% \changes{v3.0.17}{2022/02/23}{\pkg{xeCJK} 包不再自动加载 \pkg{xeCJKfntef}}
%    \begin{macrocode}
\RequirePackage{xeCJKfntef}
%    \end{macrocode}
%
% \pkg{newtx} 更新至 v1.7 版本后与 \pkg{ctex} 宏包冲突, 需要手动引入 \option{nofontspec} 选项. 判断方法为是否存在 v1.7 更新后的新宏包 \pkg{newtx.sty}
% \changes{v3.0.17}{2022/02/23}{\pkg{newtx} 更新至 v1.7 版本后与 \pkg{ctex} 宏包冲突, 需要手动引入 \option{nofontspec} 选项. 判断方法为是否存在 v1.7 更新后的新宏包 \pkg{newtx.sty}}
%    \begin{macrocode}

\RequirePackage{longtable}
\RequirePackage{booktabs}
%    \end{macrocode}
%
% 重定义 \pkg{booktabs} 宏包的默认线宽，绘制三线表时不再需要显式设置，以实现更好的内容与格式分离
% \changes{v3.0.20}{2022/05/06}{重定义 \pkg{booktabs} 宏包的默认线宽，绘制三线表时不再需要显式设置，以实现更好的内容与格式分离}
%    \begin{macrocode}
\heavyrulewidth=1.5pt
\lightrulewidth=1pt
\cmidrulewidth=0.5pt

%    \end{macrocode}
%
% \changes{v3.2a}{0000/00/00}{natbib/hyperref/url 配置从 art-deps 拆出至 art-hyperlink 模块（对齐 book-hyperlink）}
%    \begin{macrocode}
%</art-deps-a>
%    \end{macrocode}
% ^^A ======================================================================
% ^^A Module art-hyperlink: hyperref/url 配置（PDF 书签、链接颜色）（art）
% ^^A ======================================================================
%    \begin{macrocode}
%<*art-hyperlink>
\RequirePackage{hyperref}
\hypersetup{%
  CJKbookmarks=true,
  linktoc=all,
  bookmarksnumbered=true,
  bookmarksopen=true,
  bookmarksopenlevel=1,
  breaklinks=true,
  colorlinks=false,
  plainpages=false,
  pdfborder=0 0 0
}
\urlstyle{same}
%</art-hyperlink>
%<*art-deps-b>

%    \end{macrocode}
%
% \changes{v3.2a}{0000/00/00}{geometry 配置从 art-deps 拆出至 art-geometry 模块（对齐 book-geometry）}
%    \begin{macrocode}
%</art-deps-b>
%    \end{macrocode}
% ^^A ======================================================================
% ^^A Module art-geometry: 版芯尺寸与页边距（art）
% ^^A ======================================================================
%    \begin{macrocode}
%<*art-geometry>
\ExplSyntaxOn
%    \end{macrocode}
%
% \changes{v3.2a}{0000/00/00}{art-geometry 模块改用 expl3 实现。标志位仍由
%   kvoptions 声明，用 \cs{legacy_if:nTF} 桥接；深圳硕士那处的
%   \cs{ifboolexpr} 换成 \cs{bool_lazy_and:nnT}，此处不再依赖 etoolbox。}
%    \begin{macrocode}
\legacy_if:nTF { hit@debug }
  { \RequirePackage [ showframe ] { geometry } }
  { \RequirePackage { geometry } }
% 根据 PlutoThesis 原版定义而来。a4paper 即 210 * 297mm。
\geometry
  {
    a4paper , hcentering , ignoreall , nomarginpar , centering ,
    text = { 150true~mm , 236true~mm } ,
    left = 30true~mm , head = 5true~mm , headsep = 2true~mm ,
    footskip = 0true~mm , foot = 5.2true~mm
  }
\bool_lazy_and:nnT
  { \legacy_if_p:n { hit@shenzhen } }
  { \legacy_if_p:n { hit@master } }
  {
    \legacy_if:nTF { hit@opening }
      { \geometry { headsep = 25pt } }
      { \geometry { headsep = 17pt } }
  }
\ExplSyntaxOff
%</art-geometry>
%<*art-deps-c>

\ifhit@debug
  \RequirePackage{layout}
  \RequirePackage{layouts}
  \RequirePackage{lineno}
\fi

\RequirePackage{fancyhdr}
\RequirePackage{tabularx}
\RequirePackage{varwidth}
\RequirePackage{changepage}
\RequirePackage{multicol}
\RequirePackage[below]{placeins}%允许上一个section的浮动图形出现在下一个section的开始部分,还提供\FloatBarrier命令,使所有未处理的浮动图形立即被处理
\RequirePackage{flafter}       % 使得所有浮动体不能被放置在其浮动环境之前，以免浮动体在引述它的文本之前出现.
\RequirePackage{multirow}       %使用Multirow宏包，使得表格可以合并多个row格
\RequirePackage{subfigure}%支持子图 %centerlast 设置最后一行是否居中
\RequirePackage[subfigure]{ccaption} %支持双语标题
%    \end{macrocode}
% \changes{v3.0.21}{2022/05/11}{删除 \pkg{xltxtra} 宏包, 来实现规范样式的上标。\issue[125]}
% 设置深圳校区开题报告的正文行距与毕业论文模板一致
% \changes{v3.0.04}{2020/07/29}{设置深圳校区开题报告的正文行距与毕业论文模板一致}
% 此处修改与正文行距保持一致
% \changes{v3.0.05}{2020/09/13}{此处修改与正文行距保持一致}
%    \begin{macrocode}
%    \end{macrocode}
% \changes{v3.1d}{2025/03/03}{将深圳校区硕士中期报告图注实现为“图 1-1 标题”的格式}
% \changes{v3.2a}{0000/00/00}{图注/表注编号条件块从 artcls 拆出至 art-floats 模块}
%    \begin{macrocode}
%</art-deps-c>
%    \end{macrocode}
% ^^A ======================================================================
% ^^A Module art-floats: 浮动体默认行为与双语题注（art）
% ^^A ======================================================================
%    \begin{macrocode}
%<*art-floats>
% \changes{v3.2a}{0000/00/00}{art-floats 改用 expl3，两处 \cs{ifboolexpr} 换成
%   \cs{bool_lazy_and:nnT}，本模块不再依赖 etoolbox}
\ExplSyntaxOn
%    \end{macrocode}
%
% 深圳硕士中期把图注排成「图 1-1 标题」的格式。
%    \begin{macrocode}
\bool_lazy_and:nnT
  { \legacy_if_p:n { hit@shenzhen } }
  {
    \bool_lazy_and_p:nn
      { \legacy_if_p:n { hit@master } }
      { \legacy_if_p:n { hit@midterm } }
  }
  {
    \RequirePackage { caption }
    \renewcommand { \thefigure } { \thesection - \arabic { figure } }
    \captionsetup [ figure ] { labelsep=space }
  }
%    \end{macrocode}
%
% \changes{v3.2a}{0000/00/00}{schrb: 将哈尔滨校区本科开题报告图表公式编号实现为“图 X-X”、“表 X-X”、“公式 (X-X)”的格式}
%    \begin{macrocode}
\bool_lazy_and:nnT
  { \legacy_if_p:n { hit@harbin } }
  {
    \bool_lazy_and_p:nn
      { \legacy_if_p:n { hit@bachelor } }
      { \legacy_if_p:n { hit@opening } }
  }
  {
    \RequirePackage { caption }
    \renewcommand { \thefigure }   { \thesection - \arabic { figure } }
    \renewcommand { \thetable }    { \thesection - \arabic { table } }
    \renewcommand { \theequation } { \thesection - \arabic { equation } }
    \captionsetup [ figure ] { labelsep=space }
    \captionsetup [ table ]  { labelsep=space }
    \@addtoreset { figure }   { section }
    \@addtoreset { table }    { section }
    \@addtoreset { equation } { section }
  }
\ExplSyntaxOff
%</art-floats>
%<*art-deps-c>
%    \end{macrocode}
%
% 添加字号、行距等设置
%    \begin{macrocode}
%</art-deps-c>
%<*art-deps-c>
%    \end{macrocode}
%
% 更改哈尔滨校区博士硕士开题节标题为小三号字体
% \changes{v3.0.08}{2020/09/28}{更改哈尔滨校区博士硕士开题节标题为小三号字体}
%    \begin{macrocode}
%    \end{macrocode}
%
% 进行字号、行距等设置
%    \begin{macrocode}

%% 进行字号、行距等设置
%    \end{macrocode}
%
% \changes{v3.1d}{2025/03/03}{增加 \cs{parindent} 的值。\issue[201]}
% \changes{v3.2a}{0000/00/00}{章节标题格式从 artcls 拆出至 art-chapter 模块}
%    \begin{macrocode}
%</art-deps-c>
%    \end{macrocode}
% ^^A ======================================================================
% ^^A Module art-chapter: 章节标题格式（字号、缩进、间距、编号样式）（art）
% ^^A ======================================================================
%    \begin{macrocode}
%<*art-chapter>
\ExplSyntaxOn
%    \end{macrocode}
%
% \changes{v3.2a}{0000/00/00}{art-chapter 模块改用 expl3 实现。两处
%   \cs{ifboolexpr} 换成 \cs{bool_lazy_and:nnT} 与 \cs{bool_lazy_and:nnnT}，
%   本模块不再依赖 etoolbox。\cs{ctexset} 值里的 \cs{ifhit@opening} 保持原样，
%   那些条件由 ctex 延后求值，改写它们没有收益。}
%    \begin{macrocode}
\setlength { \parindent } { 2\ccwd }
\ctexset
  {
    section =
      {
        afterindent = true ,
        beforeskip = { 7.5bp } ,     % 上下空 0.5 行
        afterskip  = { 7.5bp } ,
        format = { \heiti \xiaosan [ 1.25 ] } ,
        aftername = \enspace ,
        fixskip = true ,
        break = { } ,
      } ,
    subsection =
      {
        afterindent = true ,
        beforeskip = { 7bp } , afterskip = { 7bp } ,
        format = { \heiti \sihao [ 1.25 ] } ,
        aftername = \enspace , fixskip = true , break = { } ,
      } ,
    subsubsection =
      {
        afterindent = true ,
        beforeskip = { 3bp } , afterskip = { 3bp } ,
        format = { \heiti \normalsize } ,
        aftername = \enspace , fixskip = true , break = { } ,
      } ,
    paragraph / afterindent = true ,
    subparagraph / afterindent = true
  }
%    \end{macrocode}
%
% \changes{v3.1d}{2025/03/03}{哈尔滨本科中期的一级节标题变成了中文……\issue[213]}
%    \begin{macrocode}
\bool_lazy_and:nnT
  { \legacy_if_p:n { hit@harbin } }
  {
    \bool_lazy_and_p:nn
      { \legacy_if_p:n { hit@bachelor } }
      { \legacy_if_p:n { hit@midterm } }
  }
  {
    \ctexset
      { section = { number = \chinese { section } , name = { ,、} , aftername = { } } }
  }
%    \end{macrocode}
%
% 深圳硕士的节标题格式。这一段的取值里有 \cs{ifhit@opening} 与控制空格
% \cs{\ }，控制空格在 \cs{ExplSyntaxOn} 下的记号化与平时不同，实测会让
% 变体 32 的 PDF 产生差异。所以把整块取值放在 expl3 之外定义成一个宏，
% 记号与改动前逐个相同，再由 expl3 的条件来调用。
%    \begin{macrocode}
\ExplSyntaxOff
\def\hit@art@shenzhenmaster@sectionset{%
  \ctexset{
    section={
      beforeskip=25bp,
      afterskip=27pt,
      format=\ifhit@opening\centering\xiaoer\else\xiaosan\fi\heiti,
      aftername=\ifhit@opening \quad \else \enspace \fi,
      break=\ifhit@opening \clearpage \else {} \fi
    },
    subsection={
      beforeskip=\ifhit@opening 12pt \else 15pt \fi,
      afterskip=\ifhit@opening 12pt \else 15pt \fi,
      format=\heiti\ifhit@opening \xiaosan \else \sihao \fi,
      fixskip=true,
      aftername=\ifhit@opening {\ } \else \enspace \fi
    },
    subsubsection={
      beforeskip=\ifhit@opening 12pt \else 15pt \fi,
      afterskip=\ifhit@opening 12pt \else 15pt \fi,
      format=\heiti\sihao,
      fixskip=true,
      aftername=\ifhit@opening {\ } \else \enspace \fi
    }
  }%
}
\ExplSyntaxOn
\bool_lazy_and:nnT
  { \legacy_if_p:n { hit@shenzhen } }
  { \legacy_if_p:n { hit@master } }
  { \hit@art@shenzhenmaster@sectionset }
\ExplSyntaxOff
%</art-chapter>
%    \end{macrocode}
% ^^A ======================================================================
% ^^A Module art-pagestyle: 页眉页脚样式（fancyhdr 配置）（art）
% ^^A ======================================================================
%    \begin{macrocode}
%<*art-pagestyle>
% \changes{v3.2a}{0000/00/00}{art-pagestyle 自包含 etoolbox 与 fancyhdr 依赖}
% \changes{v3.2a}{0000/00/00}{art-pagestyle 的控制流改用 expl3。页眉定义含中文
%   与控制空格，按 CONTRIBUTING §4.6 留在 expl3 之外定义成宏，再由条件调用。}
\RequirePackage{fancyhdr}
%    \end{macrocode}
% \changes{v3.1d}{2025/03/03}{为威海校区硕士开题设置章节样式与页眉页脚样式。\issue[219]}
% \changes{v3.2a}{0000/00/00}{威海硕士开题的页眉定义放入 art-pagestyle，激活仍留 artcls 保持时序}
% 两处页眉的横线画法相同，但页眉文字与适用条件不同，各自定义。
%    \begin{macrocode}
\def\hit@art@weihaimaster@pagestyle{%
  \fancypagestyle{hit@weihai@headings@main}{%
  \fancyhf{}
  \fancyfoot[C]{\xiaowu\thepage}
  \renewcommand{\headrule}{
      \vskip 1.190132pt
      \hrule\@height2.276208pt\@width\headwidth
      \vskip 0.75pt
      \hrule\@height.75pt\@width\headwidth
    }
    \fancyhead[C]{\songti\xiaowu[0]哈尔滨工业大学（威海）硕士学位论文开题报告}
    \fancyfoot[C]{-\ \thepage\ -}
  }%
}
\def\hit@art@shenzhenmaster@pagestyle{%
  \fancypagestyle{hit@shenzhen@headings@main}{%
  \fancyhf{}
  \fancyfoot[C]{\xiaowu\thepage}
  \renewcommand{\headrule}{
      \vskip 1.190132pt
      \hrule\@height2.276208pt\@width\headwidth
      \vskip 0.75pt
      \hrule\@height.75pt\@width\headwidth
    }
    \fancyhead[C]{\songti\xiaowu[0]哈尔滨工业大学（深圳）\ifhit@midterm 中期 \else 硕士学位论文开题 \fi 报告}
    \fancyfoot[C]{-\ \thepage\ -}
  }
  \fancypagestyle{hit@shenzhen@headings@front}{%
    \fancyhf{}
    \renewcommand{\headrule}{
      \vskip 1.190132pt
      \hrule\@height2.276208pt\@width\headwidth
      \vskip 0.75pt
      \hrule\@height.75pt\@width\headwidth
    }
    \fancyhead[C]{\songti\xiaowu[0]哈尔滨工业大学（深圳）硕士学位论文开题报告}
  }%
}
\ExplSyntaxOn
\bool_lazy_and:nnT
  { \legacy_if_p:n { hit@weihai } }
  {
    \bool_lazy_and_p:nn
      { \legacy_if_p:n { hit@master } }
      { \legacy_if_p:n { hit@opening } }
  }
  { \hit@art@weihaimaster@pagestyle }
\bool_lazy_and:nnT
  { \legacy_if_p:n { hit@shenzhen } }
  { \legacy_if_p:n { hit@master } }
  { \hit@art@shenzhenmaster@pagestyle }
\ExplSyntaxOff
%</art-pagestyle>
%<*art-deps-c>
\ifboolexpr{bool {hit@weihai} and bool {hit@master} and bool {hit@opening}}
  {
    \pagestyle{hit@weihai@headings@main}
  }{}
%    \end{macrocode}
% \changes{v3.1c}{2023/04/16}{为深圳校区硕士开题设置章节样式与页眉页脚样式。}
% 为深圳校区硕士开题设置章节样式与页眉页脚样式。%
% \changes{v3.1c}{2023/04/16}{为深圳校区硕士中期设置章节样式与页眉页脚样式。}
% \changes{v3.1c}{2023/04/16}{深圳校区硕士中期不换页}
% 为深圳校区硕士中期设置章节样式与页眉页脚样式。%
% \changes{v3.1d}{2025/03/03}{根据 \issue[228] 更新深圳硕士中期节标题字号和间距}
% \changes{v3.2a}{0000/00/00}{深圳硕士的 ctexset 与 fancypagestyle 被拆为两个独立条件块，分别放入 art-chapter 与 artcls，保持行为等价}
%    \begin{macrocode}
%</art-deps-c>
%<*art-deps-c>
%    \end{macrocode}

%  \changes{v3.1d}{2025/03/03}{目录格式将在下文中设置}
%    \begin{macrocode}
%     \ifboolexpr{bool {hit@midterm}}{}{\pagestyle{hit@shenzhen@headings@front}}
%    \end{macrocode}
%
% 中文封面
%    \begin{macrocode}
%</art-deps-c>
%    \end{macrocode}
% ^^A ======================================================================
% ^^A Module art-meta: 元信息字段（\hitsetup 接受的 key）（art）
% ^^A ======================================================================
%    \begin{macrocode}
%<*art-meta>
%% 中文封面
\def\hit@def@term#1{%
  \define@key{hit}{#1}{\csname #1\endcsname{##1}}
  \expandafter\gdef\csname #1\endcsname##1{%
    \expandafter\gdef\csname hit@#1\endcsname{##1}
  }
  \csname #1\endcsname{}
}

\hit@def@term{ctitlecover}
\hit@def@term{ctype}
\hit@def@term{csubject}
\hit@def@term{cauthor}
\hit@def@term{cstudentid}
\hit@def@term{cclassid}
\hit@def@term{caffil}
\hit@def@term{csupervisor}
\hit@def@term{cdate}

\def\hit@parse@keywords#1{
  \define@key{hit}{#1}{\csname #1\endcsname{##1}}
  \expandafter\gdef\csname hit@#1\endcsname{}
  \expandafter\gdef\csname #1\endcsname##1{
    \@for\reserved@a:=##1\do{
      \expandafter\ifx\csname hit@#1\endcsname\@empty\else
      \expandafter\g@addto@macro\csname hit@#1\endcsname{%
        \ignorespaces\csname hit@#1@separator\endcsname}
      \fi
      \expandafter\expandafter\expandafter\g@addto@macro%
      \expandafter\csname hit@#1\expandafter\endcsname\expandafter{\reserved@a}
    }
  }
}

\def\hitsetup{\kvsetkeys{hit}}
%</art-meta>
%    \end{macrocode}
% ^^A ======================================================================
% ^^A Module art-frontmatter: 封面与扉页（中英、本科、博士后等多版本）（art）
% ^^A ======================================================================
%    \begin{macrocode}
%<*art-frontmatter>

\newcommand{\hit@report@titlepage@graduate}{
  \ifthenelse
  {\equal{\hit@fontset}{siyuan}}%
  {\xiaosi[1]\vspace*{0.65em}}%
  {\xiaosi[1]\textcolor[rgb]{1,1,1}{\songti{\hit@hi}}}%
  \vspace*{10mm}
%    \end{macrocode}
%
% \changes{v3.0.22}{2022/05/16}{深圳校区中期模板封面范例中为“哈尔滨工业大学”。\issue[126]}
% \changes{v3.1c}{2023/04/16}{修改深圳校区硕士开题封面中的字号与字距。\issue[141]}
% \changes{v3.1c}{2023/04/16}{修改深圳校区硕士中期封面中的字号与字距。\issue[166]}
% \changes{v3.1d}{2025/03/03}{修改威海校区硕士开题封面中的字号与字距。\issue[219]}
%    \begin{macrocode}
  \begin{center}
    % \kaishu\xiaoer\textbf{\hit@cschoolname\ifhit@shenzhen\hit@shenzhencampus\fi}
    \ifboolexpr{ (bool {hit@weihai} or bool {hit@shenzhen}) and bool {hit@master} }%
    {\ziju{1/6}}{}\kaishu\xiaoer\textbf{\hit@cschoolname\ifhit@weihai （威海）\fi}
  \end{center}
  \vspace{5mm}
  \begin{center}
    \ifboolexpr{ bool {hit@shenzhen} and bool {hit@master} and bool {hit@opening} }%
    {\xiaoyi\ziju{1/8}}{\erhao}
    \songti\textbf{\hit@cxuewei\hit@cthesisname
      \ifhit@opening
        \hit@stage@opening
      \else
        \ifhit@midterm
          \hit@stage@midterm
        \fi
      \fi
      \hit@stage@doctype
    }
  \end{center}
  \vspace{10mm}
  \parbox[t][3cm][t]{\textwidth}{
    \begin{center}
      \songti\xiaoer\textbf{\hit@cthesistitleprefix\hit@title@csep\hit@ctitlecover}
      \\[12bp]
%    \end{macrocode}
%
% \changes{v3.2a}{0000/00/00}{本部硕士开题、中期封面添加“学位论文”或“实践成果”。}
%    \begin{macrocode}
     
      \ifboolexpr{ bool {hit@harbin} and bool {hit@master} and ( bool {hit@opening} or bool {hit@midterm} ) }{
        \songti\sanhao\textbf{（\hit@ctype）}
      }{}
    \end{center}
  }
  \parbox[b][3cm][t]{\textwidth}{
    \begin{center}\songti\sanhao
      \renewcommand{\arraystretch}{2.1}
      \begin{tabular}{l@{\ \ }c}
        \textbf{\hit@graduate@caffiltitle} & \underline{\makebox[6.1cm]{\textbf{\hit@caffil}}}\\
        \textbf{\hit@graduate@cmajortitle} & \underline{\makebox[6.1cm]{\textbf{\hit@csubject}}}\\
        \textbf{\hit@graduate@supervisor} & \underline{\makebox[6.1cm]{\textbf{\hit@csupervisor}}}\\
%    \end{macrocode}
%
% \changes{v3.0.03}{2020/07/03}{fix author name bug}
%    \begin{macrocode}
        \textbf{\hit@graduate@studenttitle} & \underline{\makebox[6.1cm]{\textbf{\hit@cauthor}}}\\
        \textbf{\hit@graduate@studentid} & \underline{\makebox[6.1cm]{\textbf{\hit@cstudentid}}}\\
        \textbf{\hit@graduate@datetitle} & \underline{\makebox[6.1cm]{\textbf{\hit@cdate}}}\\
      \end{tabular}\renewcommand{\arraystretch}{1}
    \end{center}
  }
  \vfill
  \ifhit@harbin
    \hit@harbin@schoolbottommark
  \else
%    \end{macrocode}
%
% \changes{v3.0.22}{2022/05/16}{深圳校区中期模板封面范例中没有“哈工大（深圳）制”字样。\issue[126]}
% \changes{v3.1c}{2023/04/16}{深圳校区硕士开题封面范例中有“深圳校区教务部~制”。}
%    \begin{macrocode}
    \ifhit@shenzhen
      \ifhit@opening
        \hit@shenzhen@master@schoolbottommark
      \fi
    \fi
  \fi
}

\newcommand{\hit@report@titlepage@bachelor}{
  \ifthenelse%
  {\equal{\hit@fontset}{siyuan}}%
  {\xiaosi[1]\vspace*{0.65em}}%
  {\xiaosi[1]\textcolor[rgb]{1,1,1}{\songti{\hit@hi}}}%
  \vspace*{10mm}
  \begin{center}
    \includegraphics[width=6.2cm]{hitlogo}
  \end{center}
%    \end{macrocode}
%
% \changes{v3.1d}{2025/03/03}{根据 \issue[213] 更改哈尔滨本科开题和中期封面，多了“本科”。}
%    \begin{macrocode}
  \begin{center}
    \songti\xiaoyi\textbf{%
      \ifhit@harbin%
        \ziju{1/12}
        本科%
      \fi%
      \hit@bachelor@cthesisname
      \ifhit@opening
        \hit@stage@opening
      \else
        \ifhit@midterm
          \hit@stage@midterm
        \fi
      \fi
      \hit@stage@doctype
    }
  \end{center}
  \vspace{15mm}
  \parbox[t][6.5cm][t]{\textwidth}{
    \begin{center}
      \songti\xiaoer\textbf{\hit@cthesistitleprefix\hit@title@csep\hit@ctitlecover}
    \end{center}
  }
  \parbox[b][6cm][t]{\textwidth}{
    \begin{center}\songti\sanhao
      \renewcommand{\arraystretch}{2.1}
%    \end{macrocode}
%
% \changes{v3.1d}{2025/03/03}{本科专业可能太长，自适应一下。\issue[199]}
%    \begin{macrocode}
      \newlength\subject@length
      \newlength\tmp@length
      \settowidth{\tmp@length}{\hit@csubject}
      \setlength{\subject@length}{\ifdim\tmp@length>6.1cm\tmp@length \else 6.1cm \fi}
    \begin{tabular}{l@{\ \ }c}
      \textbf{\hit@bachelor@cmajortitle} & \underline{\makebox[\subject@length]{\textbf{\hit@csubject}}}\\
      \textbf{\hit@bachelor@cstudenttitle} & \underline{\makebox[\subject@length]{\textbf{\hit@cauthor}}}\\
      \textbf{\hit@bachelor@cstudentidtitle} & \underline{\makebox[\subject@length]{\textbf{\hit@cstudentid}}}\\
      \ifhit@weihai % 威海校区特有
        \textbf{\hit@bachelor@cclass} & \underline{\makebox[\subject@length]{\textbf{\hit@cclassid}}}\\
      \fi
      \textbf{\hit@bachelor@csupervisortitle} & \underline{\makebox[\subject@length]{\textbf{\hit@csupervisor}}}\\
      \textbf{\hit@bachelor@cdatetitle} & \underline{\makebox[\subject@length]{\textbf{\hit@cdate}}}\\
    \end{tabular}\renewcommand{\arraystretch}{1}
  \end{center}
  }
  \vfill
  \ifhit@weihai
    \relax
  \else
    \hit@harbin@bachelor@schoolbottommark
  \fi
}

\newcommand{\hit@report@backpage@bachelor}{
  \thispagestyle{empty}
  \noindent\parbox[t][6.5cm][t]{\textwidth}{\hit@bachelor@teachercomment}
  \noindent\parbox[b][6cm][t]{\textwidth}{\hit@bachelor@teachersign\underline{\makebox[3cm]{}}\hfill\hit@bachelor@checkdate\underline{\makebox[3cm]{}}}
}

%    \end{macrocode}
%
% \changes{v3.0.22}{2022/05/16}{深圳校区中期模板目录中一级标题为黑体 \cs{heiti}，而不是加粗 \cs{bfseries}。\issue[126]}
% \changes{v3.1c}{2023/04/16}{修改错误的判断语句。\issue[141]}
% \changes{v3.1c}{2023/04/16}{深圳校区硕士开题目录有标题。\issue[166]}
% \changes{v3.1c}{2023/04/16}{深圳校区硕士开题目录修改样式。}
% \changes{v3.1c}{2023/04/16}{深圳校区硕士中期目录修改样式。\issue[166]}
% \changes{v3.1d}{2025/03/03}{威海校区硕士开题目录修改样式。\issue[219]}
% \changes{v3.2a}{0000/00/00}{目录排版从 art-frontmatter 拆出至独立模块 art-toc（对齐 book-toc）}
% 深圳校区中期模板目录中一级标题为黑体 \cs{heiti}。\par
% 深圳校区硕士开题目录有标题。
% 深圳校区硕士中期目录也有标题。
% 用 \pkg{tocloft} 宏包实现深圳校区硕士开题和中期的目录，逐步进行迁移。
%    \begin{macrocode}
%</art-frontmatter>
%<*art-toc>
\ExplSyntaxOn
%\ifboolexpr{ {bool {hit@shenzhen} or bool {hit@weihai}} and bool {hit@master} }{
  % renew \CTEX@toc@width@n,
  % add width of a space to numwidth (stored in \@tempdima)
  \cs_set_protected:Npn \CTEX@toc@width@n #1
  {
    \hbox_set:Nn \l__ctex_tmp_box {#1\,}
    \dim_set:Nn \@tempdima
      {
        \dim_max:nn { \@tempdima }
          { \box_wd:N \l__ctex_tmp_box }
      }
  }
\ExplSyntaxOff
\RequirePackage[subfigure]{tocloft}
\renewcommand{\cfttoctitlefont}{\hfill\xiaoer\heiti}
\renewcommand{\cftaftertoctitle}{\hfill\mbox{}}
\renewcommand{\cftsecfont}{\heiti}
\renewcommand{\cftsecpagefont}{}
\renewcommand{\cftsecdotsep}{\cftdotsep}
\renewcommand{\cftsecleader}{\cftdotfill{\cftsecdotsep}}
\renewcommand{\cftdotsep}{1}
\setlength{\cftaftertoctitleskip}{40pt}
\setlength{\cftbeforesecskip}{1.2ex}
\setlength{\cftbeforesubsecskip}{1.2ex}
\setlength{\cftbeforesubsubsecskip}{1.2ex}
\setlength{\cftsubsecindent}{1em}% 目录中subsection的缩进
\setlength{\cftsubsubsecindent}{2em}% 目录中subsubsection的缩进
%  \setlength{\cftsecnumwidth}{1.5em}
%  \setlength{\cftsubsecnumwidth}{0em}
%  \setlength{\cftsubsubsecnumwidth}{0em}
%    \end{macrocode}
%
% \changes{v3.1d}{2025/03/03}{根据 \issue[228] 更新深圳硕士中期目录样式}
% \changes{v3.1d}{2025/03/03}{将目录格式统一成深圳样式，除去深圳硕士开题报告的目录页有页眉}
% \changes{v3.2a}{0000/00/00}{art-toc 里选目录页页眉的 \cs{ifboolexpr} 改用
%   \cs{bool_lazy_and:nnTF}，本模块不再依赖 etoolbox}
%    \begin{macrocode}
%  \ifhit@shenzhen
%    \ifhit@midterm
\setlength{\cftsecnumwidth}{1em}% 目录中的section条目序号和标题之间的距离
\setlength{\cftsubsecnumwidth}{1.7em}% 目录中的subsection条目序号和标题之间的距离
\setlength{\cftsubsubsecnumwidth}{2.4em}% 目录中的subsubsection条目序号和标题之间的距离
%    \fi
%  \fi
%  \ifboolexpr{ bool {hit@weihai} or bool {hit@midterm}}{\tocloftpagestyle{empty}}{\tocloftpagestyle{hit@shenzhen@headings@front}}
\ExplSyntaxOn
\bool_lazy_and:nnTF
  { \legacy_if_p:n { hit@shenzhen } }
  {
    \bool_lazy_and_p:nn
      { \legacy_if_p:n { hit@master } }
      { \legacy_if_p:n { hit@opening } }
  }
  { \tocloftpagestyle { hit@shenzhen@headings@front } }
  { \tocloftpagestyle { empty } }
\ExplSyntaxOff
%}{
  % \renewcommand\tableofcontents{%
  %     \thispagestyle{empty}
  %     \ifboolexpr{ bool {hit@shenzhen} and bool {hit@master} and bool {hit@midterm} }%
  %       {\patchcmd{\l@section}{\bfseries}{\heiti}{}{}}{}
  %     \normalsize\@starttoc{toc}
  %   }
%}
%</art-toc>
%<*art-frontmatter>
%    \end{macrocode}
% \begin{macro}{\makecover}
%    \begin{macrocode}
\def\makecover{
  \begin{titlepage}
    \ifhit@bachelor
      \hit@report@titlepage@bachelor
    \else
      \hit@report@titlepage@graduate
    \fi
    \clearpage
%    \end{macrocode}
%
% \changes{v3.1c}{2023/04/16}{新增对空白页的控制。\issue[138]}
% \changes{v3.1c}{2023/04/16}{修改深圳硕士中期的页脚。\issue[166]}
%    \begin{macrocode}
    \ifhit@print \thispagestyle{empty}\mbox{} \clearpage \fi
    \ifhit@toc
      \tableofcontents
      \clearpage
    \fi
  \end{titlepage}
  \ifboolexpr{ bool {hit@shenzhen} and bool {hit@master} }{\pagestyle{hit@shenzhen@headings@main}}{}{}
}
%    \end{macrocode}
% \end{macro}
% \begin{macro}{\makebackcover}
%    \begin{macrocode}
\def\makebackcover{
  \clearpage
  \hit@report@backpage@bachelor
}
%</art-frontmatter>
%<*art-deps-c>
%    \end{macrocode}
% \end{macro}
% 引用和参考文献格式
% \changes{v3.2a}{0000/00/00}{参考文献格式从 artcls 拆出至 art-bib 模块}
%    \begin{macrocode}
%</art-deps-c>
%    \end{macrocode}
% ^^A ======================================================================
% ^^A Module art-bib: 参考文献加载与 natbib/gbt7714 设置（art）
% ^^A ======================================================================
%    \begin{macrocode}
%<*art-bib>
\RequirePackage[sort&compress,numbers]{natbib}
\RequirePackage{etoolbox}
\RequirePackage{xparse}

%% 引用和参考文献格式
\newcommand\bibstyle@numerical{\bibpunct{[}{]}{,}{s}{,}{\textsuperscript{,}}}
\newcommand\bibstyle@inline{\bibpunct{[}{]}{,}{n}{,}{\hit@inline@sep}}

\citestyle{numerical}
\DeclareRobustCommand\inlinecite{\@inlinecite}
\def\@inlinecite#1{\begingroup\citestyle{inline}\let\@cite\NAT@citenum\citep{#1}\endgroup}
\let\onlinecite\inlinecite
%    \end{macrocode}
%
% \changes{v3.1c}{2023/04/16}{深圳校区硕士开题部分参考文献没有序号。\issue[141]}
% 深圳校区硕士开题部分参考文献没有序号。其他校区的提了问题再改。
% \changes{v3.2a}{0000/00/00}{设置参考文献列表标号的对齐方式}
%    \begin{macrocode}
\ExplSyntaxOn
\renewenvironment{thebibliography}[1]{%
  \ifboolexpr{ bool {hit@shenzhen} and bool {hit@master} and bool {hit@opening} }%
  {%
    \section*{参考文献}
    \addcontentsline{toc}{section}{参考文献}
  }{}
  \list{\@biblabel{\@arabic\c@enumiv}}{%
    \str_if_eq:VnT \g_bib_labelalign_str {left} {
      \RenewDocumentCommand\makelabel{m}{##1\hfill}
    }
    \settowidth{\labelwidth}{\@biblabel{#1}}
    \setlength{\labelsep}{0.5em}
    \setlength{\itemindent}{0pt}
    \setlength{\leftmargin}{\labelsep+\labelwidth}
    \addtolength{\itemsep}{-0.8em}
    \usecounter{enumiv}%
    \let\p@enumiv\@empty
  \renewcommand\theenumiv{\@arabic\c@enumiv}}%
  \sloppy\frenchspacing
  \flushbottom
  \clubpenalty0
  \@clubpenalty \clubpenalty
  \widowpenalty0%
  \interlinepenalty-50%
  \sfcode`\.\@m
}
{\def\@noitemerr
  {\@latex@warning{Empty `thebibliography' environment}}%
\endlist\frenchspacing}
%    \end{macrocode}
% \changes{v3.2a}{0000/00/00}{控制参考文献列表中作者是否全大写和最大显示数量}
%    \begin{macrocode}
%% 用以控制参考文献列表中单个作者的格式
\exp_args:NNe \cs_new:Nn \g_bib_author_cs:n {
    \bool_if:NTF \g_bib_authorupper_bool {
        \exp_not:N \text_uppercase:n {#1}
    } {
        {#1}
    }
}
%% 用以输出参考文献列表中的作者
\NewDocumentCommand{\bibauthor}{m}{
  \int_zero:N \l_tmpa_int
  \exp_args:Ne \tl_map_inline:nn {
    \exp_args:Ne \tl_tail:n {#1}
  } {
    \int_compare:nNnT
      \l_tmpa_int = \g_bib_maxauthor_int {
      \prg_break:n {\tl_head:N {#1}}
    }
    \int_incr:N \l_tmpa_int
    \g_bib_author_cs:n {##1}
  }
  \prg_break_point:
}
\ExplSyntaxOff
%</art-bib>
%<*artcls-tail>
%    \end{macrocode}
% 杂项
% \changes{v3.2a}{0000/00/00}{misc-art 模块合并为 artcls 内联钩子}
%    \begin{macrocode}
\AtEndOfClass{\input{hithesisart.cfg}}
\AtEndOfClass{\sloppy}
%</artcls-tail>
%    \end{macrocode}
%
% \subsection{定义选项}
% \label{sec:defoption}
%    \begin{macrocode}
% %%%%%%%%%%%%%%%%%%%%%%%%%%%%%%%%%%%%%%%%%%%%%%%%%%%%%%%%%%%%%%%%%%%%%%%%%%%%%%
% bookcls
% %%%%%%%%%%%%%%%%%%%%%%%%%%%%%%%%%%%%%%%%%%%%%%%%%%%%%%%%%%%%%%%%%%%%%%%%%%%%%%
%    \end{macrocode}
% ^^A ======================================================================
% ^^A Module artcfg: art 类配置文件
% ^^A ======================================================================
%    \begin{macrocode}
%<*artcfg>
\theorembodyfont{\normalfont}
\theoremheaderfont{\normalfont\heiti}
\theoremsymbol{\ensuremath{\square}}
\newtheorem*{proof}{证明}
\theoremstyle{plain}
\theoremsymbol{}
\theoremseparator{}
\newtheorem{assumption}{假设}
\newtheorem{definition}{定义}
\newtheorem{proposition}{命题}
\newtheorem{lemma}{引理}
\newtheorem{theorem}{定理}
\newtheorem{axiom}{公理}
\newtheorem{corollary}{推论}
\newtheorem{exercise}{练习}
\newtheorem{example}{例}
\newtheorem{remark}{注释}
\newtheorem{problem}{问题}
\newtheorem{conjecture}{猜想}
\ctexset{%
  contentsname={目\hspace{\ccwd}录},
  figurename=图,
  tablename=表
}
\let\CJK@todaysave=\today
\def\CJK@todaysmall@short{\the\year 年 \the\month 月}
\def\CJK@todaysmall{\the\year 年 \the\month 月 \the\day 日}
\def\CJK@todaybig@short{\zhdigits{\the\year}年\zhnumber{\the\month}月}
\def\CJK@todaybig{\zhdigits{\the\year}年\zhnumber{\the\month}月\zhnumber{\the\day}日}
\def\CJK@today{\CJK@todaysmall}
\renewcommand\today{\CJK@today}
\newcommand\CJKtoday[1][1]{%
  \ifcase#1\def\CJK@today{\CJK@todaysave}
    \or\def\CJK@today{\CJK@todaysmall}
    \or\def\CJK@today{\CJK@todaybig}
  \fi}
%    \end{macrocode}
%
% \changes{v3.1d}{2025/03/03}{根据 \issue[220] 去掉本部答辩日期中的“日”}
%    \begin{macrocode}
\cdate{%
  \ifhit@bachelor
    \ifhit@harbin
      \CJK@todaysmall@short
    \else
      \CJK@todaysmall
    \fi
  \else \CJK@todaysmall@short \fi
}
\ifhit@doctor
\gdef\hit@cxueweishort{博}
\gdef\hit@cxuewei{\hit@cxueweishort 士}
\gdef\hit@cdegree{\hit@cxueke\hit@cxuewei}
\def\hit@cauthortitle{\hit@cxueweishort 士研究生}
\fi
\ifhit@master
\gdef\hit@cxueweishort{硕}
\gdef\hit@cxuewei{\hit@cxueweishort 士}
\gdef\hit@cdegree{\hit@cxueke\hit@cxuewei}
\def\hit@cauthortitle{\hit@cxueweishort 士研究生}
\fi
\ifhit@bachelor
\gdef\hit@cxuewei{学士}
\fi
\def\hit@stage@opening{开题}
\def\hit@stage@midterm{中期}
\def\hit@stage@doctype{报告}
%    \end{macrocode}
%
% \changes{v3.1d}{2025/03/03}{删了一些没用的定义，更改哈尔滨本科的论文名。\issue[172]}
% 咱也不知道谁想的，其他校区的本科毕设叫“毕业设计（论文）”，哈尔滨的叫“毕业论文（设计）”
%    \begin{macrocode}
\def\hit@bachelor@cxuewei{本科}
\def\hit@bachelor@cthesisname{毕业\ifhit@harbin 论文（设计）\else 设计（论文）\fi}
\def\hit@bachelor@cclass{班号}
%    \end{macrocode}
%
% \changes{v3.1d}{2025/03/03}{之前一口气删多了，恢复一下。}
%    \begin{macrocode}
\def\hit@bachelor@cmajortitle{专\hspace{2\ccwd}业}
\def\hit@bachelor@cstudenttitle{学\hspace{2\ccwd}生}
\def\hit@bachelor@cstudentidtitle{学\hspace{2\ccwd}号}
\def\hit@bachelor@caffiltitle{学院}
\def\hit@bachelor@caffiltitlesz{学院}
\def\hit@bachelor@caffiltitlewh{学院}
\def\hit@bachelor@csupervisortitle{指导教师}
\def\hit@bachelor@cdatetitle{日\hspace{2\ccwd}期}
\def\hit@bachelor@teachercomment{指导教师评语：}
\def\hit@bachelor@teachersign{指导教师签字：}
\def\hit@bachelor@checkdate{检查日期：}
\def\hit@cthesistitleprefix{题\hspace{\ccwd}目}
%    \end{macrocode}
%
% \changes{v3.1c}{2023/04/16}{修改深圳硕士开题封面“学科”。}
%    \begin{macrocode}
%    \end{macrocode}
%
% \changes{v3.2a}{0000/00/00}{修改本部硕士开题、中期封面“学院（部）”和“学科/学位专业类别”。}
%    \begin{macrocode}
\def\hit@graduate@caffiltitle{\makebox[6\ccwd][s]{\ifboolexpr{ bool {hit@harbin} and bool {hit@master} and ( bool {hit@opening} or bool {hit@midterm} ) }{学院（部}{院（系}}）}
\def\hit@graduate@cmajortitle{%
  \ifboolexpr{ bool {hit@harbin} and bool {hit@master} and ( bool {hit@opening} or bool {hit@midterm} ) }%
    {\makebox[8\ccwd][s]{学科/专业学位类别}}%
    {\makebox[6\ccwd][s]{%
      \ifboolexpr{ bool {hit@shenzhen} and bool {hit@opening} and bool {hit@master} }%
        {学科\,/\,专业}%
        {学科}%
    }}%
}
\def\hit@graduate@supervisor{\makebox[6\ccwd][s]{导师}}
\def\hit@graduate@studenttitle{\makebox[6\ccwd][s]{研究生}}
\def\hit@graduate@studentid{\makebox[6\ccwd][s]{学号}}
\def\hit@graduate@datetitle{\ifhit@opening\hit@stage@opening
\else\ifhit@midterm\hit@stage@midterm\fi\fi\hit@stage@doctype 日期}
\def\hit@graduate@enrolldate{\makebox[6\ccwd][s]{入学时间}}
\def\hit@graduate@thesistitle{\makebox[6\ccwd][s]{论文题目}}
\def\hit@graduate@cafflimajor{\makebox[6\ccwd][s]{学科专业}}
\def\hit@cschoolname{哈尔滨工业大学}
\def\hit@shenzhencampus{（深圳）}
%    \end{macrocode}
%
% \changes{v3.2a}{0000/00/00}{修改本部硕士开题封面“学位”。}
%    \begin{macrocode}
\def\hit@cthesisname{\ifboolexpr{ bool {hit@harbin} and bool {hit@opening} and bool {hit@master} }{学位}{学位论文}}
\def\hit@cthesismidtermcheck{博\hspace{.125\ccwd}士\hspace{.125\ccwd}学\hspace{.125\ccwd}位\hspace{.125\ccwd}论\hspace{.125\ccwd}文\hspace{.125\ccwd}中\hspace{.125\ccwd}期\hspace{.125\ccwd}检\hspace{.125\ccwd}查\hspace{.125\ccwd}报\hspace{.125\ccwd}告}
\def\hit@weihaicampus{（威海）}
\def\hit@cschoolnametitle{\ifboolexpr{bool {hit@harbin} and bool {hit@bachelor}}{学校}{授予学位单位}}
\def\hit@cdatetitle{答辩日期}
\def\hit@caffiltitle{所在单位}
\def\hit@csubjecttitle{学科}
\def\hit@cdegreetitle{申请学位}
\def\hit@csupervisortitle{导师}
\def\hit@cassosupervisortitle{副导师}
\def\hit@ccosupervisortitle{联合导师}
\def\hit@title@csep{：}
\def\hit@natclassifiedindextitle{国内图书分类号}
\def\hit@internatclassifiedindextitle{国际图书分类号}
\def\hit@secretlevel{密级}
\def\hit@schoolidtitle{学校代码}
\def\hit@schoolid{10213}
\newcommand{\hit@authorsig}{作者签名：}
\newcommand{\hit@teachersig}{导师签名：}
\newcommand{\hit@frontdate}{日期：}
\newcommand{\hit@authorizationtitle}{学位论文使用权限}
%    \end{macrocode}
%
% \changes{v3.0.22}{2022/05/16}{深圳校区中期模板封面范例中没有“哈工大（深圳）制”字样。\issue[126]}
%    \begin{macrocode}
% \newcommand{\hit@shenzhen@schoolbottommark}{
%    \begin{center}
%      \songti\sanhao\textbf{哈工大（深圳）制}
%    \end{center}
%    \begin{center}
%      \songti\sanhao\textbf{二〇一二年三月}
%    \end{center}
% }
%    \end{macrocode}
%
% 此处删除了时间
% \changes{v3.0.06}{2020/09/13}{此处删除了时间}
%    \begin{macrocode}
\newcommand{\hit@harbin@schoolbottommark}{
   \begin{center}
     \songti\sanhao\textbf{研究生院制}
   \end{center}
}
%    \end{macrocode}
%
% \changes{v3.1c}{2023/04/16}{根据 \issue[151] 将本部本科开题下方的标识更换为 \texttt{.eps} 文件。}
%    \begin{macrocode}
\newcommand{\hit@harbin@bachelor@schoolbottommark}{
   \begin{center}
     % \lishu\xiaoer\textbf{哈尔滨工业大学教务处制}
     \includegraphics{hrb-bachelor-bottommark.eps}
   \end{center}
}
%    \end{macrocode}
%
% \changes{v3.1c}{2023/04/16}{深圳校区硕士开题模板含有“深圳校区教务处 制”。\issue[141]}
%    \begin{macrocode}
\newcommand{\hit@shenzhen@master@schoolbottommark}{
   \begin{center}
     \songti\sanhao\textbf{深圳校区教务部~制}
   \end{center}
}
\newcommand{\hit@shenzhen@doctor@midterm@note}{
\thispagestyle{empty}
{\begin{center}
\songti\sanhao\textbf{说\hspace{4\ccwd}明}
\end{center}
}
\vspace{1cm}\songti\sihao[2]
\begin{enumerate}[itemsep=-4bp]
\item 博士研究生学位论文中期报告一般在入学后的第五学期末完成。
\item 此报告由学生本人填写，请导师签字后提交一份给检查小组。答辩结束后由检查组长签署意见。
\item 报告经中期报告检查组长签字同意后，由学科部保存，以备论文答辩时参考。研究生院学生培养处将对研究生的学位论文中期报告进行抽查。
\item 报告统一用A4纸打印。
\end{enumerate}
}
\newcommand{\hit@datefill}{\hspace{2.5em}年\hspace{1.5em}月\hspace{1.5em}日}
\def\hit@hi{嗨！thesis}
\def\hithesis{\textsc{hi}\-\textsc{Thesis}}
\def\hit{哈尔滨工业大学}
%</artcfg>
%    \end{macrocode}
